\chapter{Langues peu dotées}
une langue peu dotée est une langue 
disposant de peu de ressources linguistiques 
informatisées. (définition wiktionnaire)

Les langues peu dotées sont des langues pour lesquels il est difficile de trouver des ressources informatisées.
Ces langues, sont au centre de la recherche linguistique mené dans ce mémoire. En effet, ces langues 
(ici Corse, Occitans et Alsacien) sont des langues qu'il est intéressant d'étudier. Ces langues posent diverses
problématiques qui, si résolue, Les langues peu dotées n'ont pas fait 
permettrait de faire avancer la recherche sur le traitement automatique des
langues. Les problématiques en question ont déjà était discuter par l'état de l'art.

\section{Problématiques}
Comme l'on précisé par le passé Rapp, Sharoff et Zweigenbaum (2016) : Parmis les environs 7000 langues, 
des quelles 600 ont une forme écrite, la vaste majorité sont du type "ressources basses".


%todo \/
% petites tailles de corpus
% grammaire différente pour un mm mot.